\documentclass[11pt]{article}
\usepackage[a4paper,margin=2.2cm]{geometry}
\usepackage{amsmath,amssymb,booktabs}
\usepackage{algorithm,algorithmic}
\usepackage[hidelinks]{hyperref}
\setlength{\parindent}{0pt}
\setlength{\parskip}{0.6em}
\title{PMSM Project: Week 4 Mathematical Review}
\author{AI-assisted mathematical appendix draft}
\date{2 October 2026}
\begin{document}
\maketitle
\textbf{Scope.} This standalone appendix is a source for the team's report and viva.
It uses documented teaching parameters. It is not the complete team report or a
statement of any student's personal contribution. All final identities, parameters,
Git evidence and submission records must be checked by the team.

\section{Model, initial data and reference solution}
The state is $y=(i_d,i_q)^T$. Electrical angular speed $\omega_e$ is supplied externally.
With constant voltage and coefficients, the current subsystem is
\begin{align}
L_d\dot i_d&=u_d-Ri_d+\omega_eL_qi_q,\\
L_q\dot i_q&=u_q-Ri_q-\omega_e L_di_d-\omega_e\psi_f.
\end{align}
Use A, V, $\Omega$, H, Wb and seconds. Electrical speed satisfies
$\omega_e=p\omega_m$; mechanical rpm cannot be substituted directly.
The affine IVP is $\dot y=Ay+b$, $y(0)=y_0$, where
\[
A=\begin{pmatrix}-R/L_d&\omega_eL_q/L_d\\-\omega_eL_d/L_q&-R/L_q\end{pmatrix},
\qquad b=\begin{pmatrix}u_d/L_d\\(u_q-\omega_e\psi_f)/L_q\end{pmatrix}.
\]
For positive $R,L_d,L_q$, $f(t,y)=Ay+b$ is globally Lipschitz and the IVP has a
unique solution for all finite times. Solve $Ay^*=-b$ without forming a numerical inverse.
The exact affine-system formula is
\[
y(t)=y^*+\exp(At)(y_0-y^*).
\]
The teaching case has $R=0.5$, $L_d=0.003$, $L_q=0.004$, $\psi_f=0.05$,
$\omega_e=1000$, $u_d=0$, $u_q=60$, $y_0=0$ and $T=0.03$ in SI units.
The equilibrium is $(3.2653061224,0.4081632653)^T$ A. A tight Radau trajectory and
the zero-speed decoupled exponential special case provide independent checks.

\newpage
\section{Continuous perturbation stability}
The trace and determinant satisfy
\[
\operatorname{tr}A=-R(1/L_d+1/L_q)<0,\qquad
\det A=R^2/(L_dL_q)+\omega_e^2>0.
\]
Thus the characteristic roots are
\[
\lambda_\pm=-\frac R2(1/L_d+1/L_q)\pm
\sqrt{\frac{R^2}{4}(1/L_d-1/L_q)^2-\omega_e^2}.
\]
For the teaching case these are $-145.8333333\pm999.7829626\,\mathrm{i}$
$\mathrm{s}^{-1}$. Their equal decay rates do not show a large separation of decay scales.
The explicit-step restriction can instead be driven by oscillation.

Set $v=y-y^*$, $D=\operatorname{diag}(L_d,L_q)$ and $x=Dv$. Then
\[
\dot x=Bx,\qquad B=DAD^{-1}=
\begin{pmatrix}-R/L_d&\omega_e\\-\omega_e&-R/L_q\end{pmatrix}.
\]
With $V=\tfrac12x^Tx$, the skew coupling cancels:
\begin{align}
\dot V&=-\frac R{L_d}x_d^2-\frac R{L_q}x_q^2\\
&=-R(L_dv_d^2+L_qv_q^2)
\le-2\min(R/L_d,R/L_q)V.
\end{align}
Consequently $V(t)\le V(0)\exp[-2\min(R/L_d,R/L_q)t]$.
The equivalent checked identity, for $P=D^2$, is
\[
A^TP+PA=-2R\operatorname{diag}(L_d,L_q).
\]
This is a global result for the stated constant-coefficient perturbation system.
It requires a fresh derivation for variable speed, state-dependent inductance or
mechanical feedback. The quadratic function is a weighted perturbation measure,
not the energy of a complete electrical-mechanical motor.

\section{One-step defect and global error}
Taylor substitution into explicit Euler gives a one-step defect
$\tfrac12h^2 y''+O(h^3)$. At fixed $T$, a Lipschitz error recursion and discrete
Gronwall yield global order one. If LTE is defined after dividing the defect by $h$,
its order is one; state the convention explicitly.
Classical RK4 has one-step defect $O(h^5)$ and global order four under smoothness
and stability assumptions. For $v'=Av$, its update is the fourth-degree Taylor
polynomial of $\exp(hA)$; four function calls alone do not establish fourth order.

\newpage
\section{Absolute stability and the measured boundary}
For $u'=\lambda u$, $z=h\lambda$, the stability functions are
\begin{align}
\mathcal R_E(z)&=1+z,\\
\mathcal R_{RK4}(z)&=1+z+z^2/2+z^3/6+z^4/24,\\
\mathcal R_{IE}(z)&=(1-z)^{-1}.
\end{align}
The sets are $|1+z|\le1$, $|\mathcal R_{RK4}(z)|\le1$ and $|1-z|\ge1$,
respectively. The implicit pole $z=1$ is excluded. Implicit Euler contains the
left half-plane and damps $|z|\to\infty$ there, but a stable step can remain inaccurate.

For $\lambda=-\alpha+\mathrm{i}\beta$, $\alpha>0$,
\[
|1+h\lambda|^2=(1-\alpha h)^2+\beta^2h^2\le1
\quad\Longrightarrow\quad 0<h\le\frac{2\alpha}{\alpha^2+\beta^2}.
\]
The teaching case has $h_E=0.2857142857$ ms; RK4 first exits its stability set
on the same eigenvalue ray at $2.9296218978$ ms. Equality need not imply decay.
The saved experiment tests $0.8$ and $1.2$ times each explicit boundary.
For general nonlinear systems, an $h\lambda(J_f)$ overlay is a local frozen-Jacobian
diagnostic. The stronger constant-matrix statement here concerns powers of
$\mathcal R(hA)$, not a proof for arbitrary nonlinear trajectories.

\section{Newton residual and local adaptivity}
For implicit Euler the new state is a root of
\[
F(w)=w-y_n-hf(t_n+h,w),\qquad J_F=I-hJ_f(t_n+h,w).
\]
Solve $J_F\Delta=-F$, then update $w\leftarrow w+\Delta$. The implemented residual
tolerance is $10^{-12}$ in current units. This affine model needs one exact-arithmetic
Newton correction; that observation does not test nonlinear quadratic convergence.

With one coarse step and two half steps, the fine-step error estimate is
$(y_f-y_c)/(2^p-1)$. Here $p=1$, so the denominator is one. The acceptance ratio is
\[
E=\max_j\frac{|y_{f,j}-y_{c,j}|}
{\mathrm{atol}_j+\mathrm{rtol}\max(|y_{n,j}|,|y_{f,j}|)}.
\]
Only $E\le1$ accepts the two-half-step state. Rejection preserves the old time
and state. The controller uses safety factor $0.9$, exponent $-1/2$, and bounded
step ratios. Local control is not a bound on accumulated global error.

\newpage
\section{Reproduction and measured cost}
The recorded endpoint error is
\[
E(h)=\max_{j\in\{d,q\}}|i_{j,h}(T)-i_{j,\mathrm{ref}}(T)|.
\]
The fit of $\log E$ to $\log h$ over $120\le N\le7680$ gives
$4.0375549205$ with slope standard error $0.0090768069$. This is $4.04$ to three
significant digits, and $4.038\pm0.009$ in the figure's decimal display.
The very finest points are excluded because floating-point error becomes non-monotone.
The regression band is a deterministic fit diagnostic, not repeated-sample uncertainty.

\texttt{python code/week4\_reproduction.py} recomputes the values, compares them
with \texttt{report/report\_claims.json}, and writes the matching lines.
Its strict mode checks a real expected Git commit, source tracking and a clean
folder before running. A content hash complements the commit; it does not replace it.

The matched-accuracy study imposes shared maximum sampled error ceilings against
tight Radau ($\mathrm{rtol}=10^{-12}$, $\mathrm{atol}=10^{-14}$).
The first qualifying grid in a powers-of-two refinement sequence is selected.
The measured values for the teaching case are:
\begin{center}
\begin{tabular}{llrrr}\toprule
Target [A]&Method&$N$&Error [A]&RHS calls\\\midrule
0.01&Euler&15360&0.00827688&15360\\
0.01&RK4&60&0.00473035&240\\
0.01&Implicit Euler&15360&0.00822511&30720\\
0.005&Euler&30720&0.00413195&30720\\
0.005&RK4&60&0.00473035&240\\
0.005&Implicit Euler&30720&0.00411900&61440\\\bottomrule
\end{tabular}
\end{center}
RHS counts include Newton residual evaluations. Jacobian calls and Newton updates
are reported separately. Linear solves, reference construction, setup and plotting
are excluded; no wall-clock ranking is claimed. These are shared error ceilings,
not exactly equal attained errors. Errors are measured on each solver grid,
not proven maxima between nodes. Recompute this table if final parameters change.

\section{References and report integration}
The mathematical conventions follow the supplied Lecture 1 and Lecture 2 materials
and the Project 1 pseudocode handout. Reproduction and final verification follow
Tutorial 4, the Project Brief and the Submission Checklist. General numerical
analysis background is available in D. F. Griffiths and D. J. Higham,
\emph{Numerical Methods for Ordinary Differential Equations: Initial Value Problems},
Springer, 2010 (provided course copy).
The final team report must cite the actual editions/materials used, attach a team
AI log and one labelled, genuine ICS per member, and satisfy the minimum eight-page
requirement. This mathematical appendix is only part of that report.

\newpage
\section{Pseudocode matching the implementation}
The following uses the \texttt{algorithmic} package, assignment arrows, mathematical
subscripts and one statement per line. The last step is shortened to reach $T$.
\begin{algorithm}[H]
\caption{Explicit Euler with endpoint handling}
\begin{algorithmic}[1]
\REQUIRE $f,y_0,t_0,T,h$ with $t_0<T$, $h>0$
\STATE $t\leftarrow t_0$
\STATE $y\leftarrow y_0$
\WHILE{$t<T$}
\STATE $s\leftarrow\min(h,T-t)$
\STATE $y\leftarrow y+s f(t,y)$
\STATE $t\leftarrow t+s$
\ENDWHILE
\STATE \textbf{return} $y$
\end{algorithmic}
\end{algorithm}
\begin{algorithm}[H]
\caption{One classical RK4 step}
\begin{algorithmic}[1]
\REQUIRE $f,t,y,h$
\STATE $k_1\leftarrow f(t,y)$
\STATE $k_2\leftarrow f(t+h/2,y+hk_1/2)$
\STATE $k_3\leftarrow f(t+h/2,y+hk_2/2)$
\STATE $k_4\leftarrow f(t+h,y+hk_3)$
\STATE $w\leftarrow y+h(k_1+2k_2+2k_3+k_4)/6$
\STATE \textbf{return} $w$
\end{algorithmic}
\end{algorithm}
\begin{algorithm}[H]
\caption{Implicit Euler via Newton; residual tolerance $10^{-12}$ A}
\begin{algorithmic}[1]
\REQUIRE $f,J_f,t,y,h,\mathrm{tol},\mathrm{maxiter}$
\STATE $w\leftarrow y$
\FOR{$k=0,\ldots,\mathrm{maxiter}$}
\STATE $F\leftarrow w-y-h f(t+h,w)$
\IF{$\|F\|_\infty<\mathrm{tol}$}
\STATE \textbf{return} $w$
\ENDIF
\IF{$k=\mathrm{maxiter}$}
\STATE \textbf{raise} Newton failure
\ENDIF
\STATE $J_F\leftarrow I-hJ_f(t+h,w)$
\STATE $\Delta\leftarrow\operatorname{solve}(J_F,-F)$
\STATE $w\leftarrow w+\Delta$
\ENDFOR
\end{algorithmic}
\end{algorithm}

\newpage
\begin{algorithm}[H]
\caption{Adaptive implicit Euler; accept two half steps, $p=1$}
\begin{algorithmic}[1]
\REQUIRE $f,J_f,t_0,T,y_0,h_0,\mathrm{atol},\mathrm{rtol}$
\STATE $t\leftarrow t_0$
\STATE $y\leftarrow y_0$
\STATE $h\leftarrow h_0$
\WHILE{$t<T$}
\STATE $h\leftarrow\min(h,T-t)$
\STATE $y_c\leftarrow\operatorname{ImplicitStep}(t,y,h)$
\STATE $y_m\leftarrow\operatorname{ImplicitStep}(t,y,h/2)$
\STATE $y_f\leftarrow\operatorname{ImplicitStep}(t+h/2,y_m,h/2)$
\STATE $s_j\leftarrow\mathrm{atol}_j+\mathrm{rtol}\max(|y_j|,|y_{f,j}|)$
\STATE $E\leftarrow\max_j|y_{f,j}-y_{c,j}|/s_j$
\IF{$E\le1$}
\STATE $t\leftarrow t+h$
\STATE $y\leftarrow y_f$
\ENDIF
\IF{$E=0$}
\STATE $q\leftarrow2$
\ELSE
\STATE $q\leftarrow\min(2,\max(0.2,0.9E^{-1/2}))$
\ENDIF
\IF{$E>1$}
\STATE $q\leftarrow\min(0.5,q)$
\ENDIF
\STATE $h\leftarrow hq$
\ENDWHILE
\STATE \textbf{return} accepted times and states
\end{algorithmic}
\end{algorithm}
On Newton failure, the implementation records the failed trial, leaves $t,y$
unchanged, and shrinks $h$ with the minimum rejection factor. It also enforces
a minimum step and a maximum trial count. The pseudocode displays the successful
trial path; these safeguards remain part of the runnable implementation.

\textbf{Final viva check.} Derive one stability bound without code; explain the
state units; distinguish algebraic, local and global errors; reproduce a reported
number from the agreed commit; and identify the costs excluded from the efficiency
comparison. Each student records their own work and evidence.
\end{document}
